\section{Polyglot Discipline}
\label{sec:polyglot}

Every \pqtoolkit{} project ships byte-identical reference implementations in eight languages, all consumed by the same cross-language KAT harness. This section documents the discipline that makes that work.

\subsection{The authority order}

When any two sources disagree, conflicts are resolved in the following order:

\begin{enumerate}
\item The project's whitepaper LaTeX (this document for cross-cutting concerns; the per-project document for project-specific normative content).
\item The project's Python reference: \texttt{pqid/python/pqid.py}, \texttt{pqhash/python/pqhash.py}, \texttt{pqcrypt/python/pqcrypt.py}.
\item The project's canonical C ABI header: \texttt{pqid/c/pqid.h}, \texttt{pqhash/c/pqhash.h}, \texttt{pqcrypt/c/pqcrypt.h}.
\item Every other language port.
\end{enumerate}

A cross-language disagreement is never resolved by changing the spec to match an outlier. The outlier is fixed. ``Python is right'' is the short version of the rule, with whitepaper-overrides-Python and C-overrides-other-languages as the tie-breakers.

\subsection{The eight languages}

Every project ships reference implementations in:

\begin{center}
\small
\begin{tabular}{@{}lll@{}}
\toprule
\textbf{Language} & \textbf{Role} & \textbf{Why} \\
\midrule
Python      & Readable canonical reference & The spec is the same shape as the Python; the Python is the test. \\
C           & Canonical ABI                 & Offsets, error codes, function signatures pinned in \texttt{pqcrypt.h} / \texttt{pqid.h} / \texttt{pqhash.h}. \\
C++         & C-library wrapper             & RAII surface over C; no normative content of its own. \\
Go          & Stdlib-heavy port             & Tests that the spec is implementable using only stdlib + audited deps. \\
Rust        & Strongest constant-time toolchain & \code{\#![forbid(unsafe\_code)]}; tests the safety story. \\
TypeScript  & Browser + Node                 & Tests that the spec is implementable in JavaScript runtimes. \\
Java        & JVM                            & Tests that the spec works against JDK stdlib (SunJCE) where applicable. \\
CPython extension & Optimized Python      & Mirrors the Python API one-to-one; provides a C-fast option without forcing C bindings. \\
\bottomrule
\end{tabular}
\end{center}

For \pqcrypt{} specifically (post-quantum, more moving parts than the other two projects), the alpha language set during v1.x development was Python + C + Go + Rust --- the minimum that exercises a readable reference, a canonical ABI, a stdlib-heavy ecosystem, and the strongest constant-time toolchain. C++/TypeScript/Java/CPython ext joined in v1.1.

\subsection{The cross-language KAT harness}

\texttt{tools/conformance/run.sh} is the single entry point. It discovers per-language runners under \texttt{tools/conformance/python/}, \texttt{tools/conformance/c/}, etc., builds each one if needed, and runs each against every KAT file in \texttt{pqid/conformance/vectors/}, \texttt{pqhash/conformance/vectors/}, and \texttt{pqcrypt/conformance/vectors/}.

\begin{lstlisting}[language=bash]
# Full run: all eight languages times all three projects
bash tools/conformance/run.sh

# Subset of languages
bash tools/conformance/run.sh python rust go

# Skip languages whose toolchain isn't available locally
PQTK_CONFORMANCE_SKIP="cpython java" bash tools/conformance/run.sh

# PQCrypt v1.x alpha quick loop
PQTK_CONFORMANCE_SKIP="cpython java cpp typescript" bash tools/conformance/run.sh
\end{lstlisting}

\subsection{Branch protection}

\texttt{.github/workflows/conformance.yml} runs the harness on every PR against \texttt{main}. Red harness $\Rightarrow$ no merge. Per-language workflows (\texttt{python.yml}, \texttt{c.yml}, \ldots) cover per-OS and per-toolchain variants independently; the cross-language harness is the only gate that catches polyglot drift.

PQCrypt has an additional per-project workflow (\texttt{pqcrypt.yml}, added in PQCK-00023) that runs the PQCrypt-only matrix across Ubuntu and macOS for the alpha languages plus the LaTeX whitepaper build.

\subsection{What a spec change looks like}

A normative change to any project follows this order:

\begin{enumerate}
\item Update the project's LaTeX whitepaper (and, if cross-cutting, this umbrella document).
\item Update the Python reference.
\item Update the C reference (header first if offsets/constants change, then \texttt{.c}).
\item Port to every other language in lockstep. Each must produce byte-identical output for the same input as Python.
\item Regenerate conformance vectors from the Python reference; commit verbatim.
\item Run \texttt{bash tools/conformance/run.sh}; confirm green.
\item Write a release-notes entry under \texttt{<project>/docs/release-notes/YYYY-MM-DD-<slug>.md}; append to \texttt{<project>/docs/release-notes/CHANGELOG.md}.
\item For security-relevant changes, tag the PR \texttt{sec(<TICKET>)} and add a ``Security considerations'' section to the release-notes entry.
\end{enumerate}

Optimization PRs that only touch one language are fine and don't require porting, as long as observable behavior is identical and KAT vectors still pass byte-identical.
